% !TEX root = ../main/main.tex
\chapter{Conclusion and Future Work}
\label{chap:conclusion}

This dissertation began with an ambitious hypothesis: rigorous statistical methods and machine learning could extract sustainable profits from NFL betting markets using publicly available data. After 5,529 games, 21 seasons, and a 13-month research journey through 11 model configurations, we arrive at a nuanced but ultimately \textbf{negative conclusion for spread betting profitability}---though with significant methodological contributions that advance the field.

\textbf{Critical finding (December 2025):} The cleaned baseline from the comprehensive data leakage investigation (\S\ref{sec:data-leakage-investigation}) shows \textbf{51.0\% win rate and $-7.5$\% ROI}---below the 52.4\% breakeven threshold. Earlier claims of +2.96\% ROI were computed before leakage remediation and are invalidated. This dissertation confirms \textbf{semi-strong market efficiency} for NFL spread betting with publicly available data.

\section{What We Learned}

\subsection{The Central Lesson: Market Efficiency Confirmed Through Rigorous Validation}

The journey from v1.0 to v6.0, culminating in the comprehensive data leakage investigation, demonstrates that NFL spread betting markets are \textbf{efficient for publicly available information}. Our initial v1.0 system suffered catastrophic failure (maximum drawdown = $-83.5\%$). Subsequent iterations appeared to show improvement, but the final leakage investigation revealed that earlier positive results were artifacts of data contamination.

\textbf{The cleaned baseline shows:}
\begin{itemize}
  \item \textbf{Win rate:} 51.0\% (below 52.4\% breakeven threshold)
  \item \textbf{ROI:} $-7.5\%$ (confirming negative expected value after vig)
  \item \textbf{True out-of-sample (2025 holdout):} 74.9\% spread \emph{prediction} accuracy---but prediction accuracy does not imply betting profitability
  \item \textbf{Props performance:} 44.25 MAE with equal-weight averaging (props market may be less efficient, but pipeline is currently non-functional---model files missing, data tables empty)
  \item \textbf{Complete 2025 season (February 2026):} Full evaluation across 285 games (272 regular season + 13 postseason including Super Bowl~LIX) confirms the market efficiency finding with $-3.2\%$ ATS ROI on 208 qualified bets (\S\ref{subsec:postseason-validation})
\end{itemize}

Earlier claims of +2.96\% ROI and 53.89\% win rate were computed before comprehensive leakage remediation and are now \textbf{INVALIDATED}. The Ridge Residuals method---initially appearing to offer 88\% improvement---was revealed to contain target leakage that inflated metrics by $>$14$\times$.

\textbf{The key insight}: Market efficiency is real. Publicly available NFL data---play-by-play, injuries, weather, historical odds---is insufficient to overcome the 4.5\% vigorish barrier. This is not a failure of methodology but a \emph{confirmation of efficient market theory}.

\textbf{The implication}: The dissertation's contribution is \textbf{methodological}, not economic. We demonstrate (1) how to rigorously validate ML systems in adversarial environments, (2) how to detect and remediate data leakage that persists across multi-season backtests, and (3) how to honestly report negative results. These contributions advance the field even though profitability was not achieved.

\subsection{Methodological Contributions}

Beyond profitability, this dissertation makes four methodological contributions that generalize beyond NFL betting:

\begin{center}
\fbox{\begin{minipage}{0.95\textwidth}
\vspace{6pt}
\textbf{Core Contributions (echoed from Chapter 1)}
\begin{enumerate}
  \item \textbf{Rigorous negative results}: Weather has no predictive value; calibration does not imply profitability; RL provides marginal gains over simpler Kelly baselines (\Cref{chap:results}, Table~\ref{tab:weather-effects}, Table~\ref{tab:rl-vs-baseline}).

  \item \textbf{Complete system architecture}: A full pipeline from ingestion to evaluation (\Cref{chap:data}, \Cref{chap:sim}), reusable for alternative domains (player props, lower-vig exchanges).

  \item \textbf{Dependence-aware evaluation}: Copula-based methods for correlated outcomes (\Cref{chap:risk}, \S\ref{subsec:teaser-copula}) demonstrate proper modeling of same-game parlays and teasers.

  \item \textbf{Transparent failure analysis}: Documentation of when and why the system declines to act (\Cref{chap:results}, Table~\ref{tab:zero-weeks}, 21\% zero-bet weeks) prevents overfitting to lucky backtests.
\end{enumerate}
\vspace{6pt}
\end{minipage}}
\end{center}

\paragraph{1. Rigorous Negative Results (Expanded).}
We transparently document three significant null findings:
\begin{itemize}
  \item \textbf{Weather has no predictive value}: Comprehensive analysis of 1,021 outdoor games shows wind ($r=0.004$, $p=0.90$) and temperature ($r=0.055$, $p=0.08$) have no significant correlation with scoring. Modern NFL teams neutralize weather effects through preparation, and betting markets efficiently price residual impacts.
  \item \textbf{Calibration does not imply profitability}: Brier=0.2515 and CLV=+14.9bps are insufficient to overcome vig. This clarifies the gap between statistical performance and economic viability.
  \item \textbf{RL provides marginal gains}: Reinforcement learning improved Sharpe by $\sim$0.1--0.2 over Kelly baselines but required 10--40 hours of compute per training run. Given modest gains, simpler Kelly-LCB baselines are preferred for production.
\end{itemize}

These negative results prevent future researchers from wasting effort on low-value features and overstated RL claims.

\paragraph{2. Complete Betting System Architecture.}
We demonstrate a full pipeline from data ingestion (TimescaleDB) to model training (GLM/XGBoost/state-space) to risk management (CVaR LP, Kelly sizing) to evaluation (OPE, simulator acceptance). This infrastructure can be reused for alternative domains (player props, portfolio optimization) even if NFL profitability remains elusive.

\paragraph{3. Dependence-Aware Evaluation.}
Our copula-based approach to same-game parlays and teasers (\Cref{chap:risk}) demonstrates how to model correlated outcomes properly. Ignoring dependence overstates teaser EV by 2--5 percentage points—a critical correction for multi-leg betting strategies.

\paragraph{4. Transparent Failure Analysis.}
We document when and why the system fails: 21\% of weeks produce zero bets due to OPE gating (conservative lower bounds $\le 0$) or simulator rejection (CVaR/drawdown breach). This transparency prevents overfitting to lucky backtests and enforces operational discipline.

\section{Limitations and Threats to Validity}

\subsection{Data Limitations}
\begin{itemize}
  \item \textbf{No proprietary tracking data}: Our models use publicly available play-by-play, injuries, and weather. Access to player tracking (Next Gen Stats), formation data, or insider injury intel would improve edge.
  \item \textbf{Historical odds survivorship}: Delisted sportsbooks create potential selection bias. We retain NA fills but cannot fully mitigate this.
  \item \textbf{Retrospective injury revisions}: Official injury reports sometimes update retroactively, creating look-ahead bias. We use as-of snapshots but acknowledge residual risk.
\end{itemize}

\subsection{Model Limitations}
\begin{itemize}
  \item \textbf{Linear correlation assumptions}: Gaussian and t-copulas capture tail dependence but may miss complex non-linear dependencies in extreme scenarios.
  \item \textbf{State-space over-smoothing}: Bayesian state-space models excel at temporal stability but may lag rapid regime shifts (coaching changes, injury clusters).
  \item \textbf{RL sample efficiency}: Offline RL struggles with small sample sizes. NFL's 272 games/season limits training data compared to high-frequency domains.
\end{itemize}

\subsection{External Validity}
\begin{itemize}
  \item \textbf{Market regime shifts}: Results generalize across 2004--2024 but may not hold if betting markets fundamentally change (regulation shifts, algorithm-driven pricing).
  \item \textbf{Execution assumptions}: We assume fill rates and slippage based on historical patterns. Live execution may deviate during volatile periods (line steam, limit reductions).
  \item \textbf{Simulated liquidity}: Acceptance tests use pessimistic friction regimes but cannot perfectly replicate real-world market conditions.
\end{itemize}

\subsection*{Limitations and Boundaries of the Playbook}
\addcontentsline{toc}{subsection}{Limitations and Boundaries of the Playbook}

While this dissertation presents a generalizable "playbook" for building profitable AI in competitive environments, its success is contingent on specific characteristics of the target domain. The principles of feature engineering, ensemble diversification, and risk management are universal, but their application must be adapted. The NFL betting market proved an ideal testbed for several reasons:

\begin{itemize}
    \item \textbf{Rich, Structured Public Data:} The availability of decades of high-quality play-by-play data is a prerequisite. Domains with sparse or unstructured data would require significantly more investment in data processing.
    \item \textbf{Medium-Frequency Dynamics:} Game outcomes are weekly events. This cadence allows for complex models to be trained and validated offline, a luxury not available in high-frequency domains like algorithmic trading, where microsecond latency and continuous model updates are paramount.
    \item \textbf{A "Solvable" Level of Efficiency:} The NFL market is highly efficient, but the 4.5\% vigorish creates a fixed barrier that can be overcome with a superior model. In markets with near-zero transaction costs and adversarial participants (e.g., quantitative finance), the alpha from such a system would likely decay much faster.
\end{itemize}

Applying this playbook to other domains would require careful consideration of these factors. In supply chain optimization, for instance, the challenge may not be market efficiency but modeling real-world physical constraints and stochastic delays. In finance, the emphasis would shift from complex feature engineering to low-latency infrastructure and adversarial game theory. The core lesson holds—know your domain, diversify your models, manage your risk—but the implementation details will differ significantly.

\section{Future Directions}

Despite current unprofitability, this work opens several research directions:

\subsection{Alternative Markets}
\begin{itemize}
  \item \textbf{Lower-vig exchanges}: Test methods on betting exchanges (Betfair, Pinnacle) where vig is 1--2\% instead of 4.5\%. Our CLV=+14.9bps may suffice at lower friction.
  \item \textbf{Player props and derivatives}: Apply copula methods to correlated player performance (QB passing yards + receiver yards). These markets may be less efficient than game-level spreads.
  \item \textbf{Live in-game betting}: Extend RL to dynamic in-game markets where information arrives continuously (score updates, injury substitutions).
\end{itemize}

\subsection{Methodological Extensions}
\begin{itemize}
  \item \textbf{Multi-league transfer learning}: Train models on NBA/MLB data and transfer features to NFL. Cross-league learning may improve sample efficiency.
  \item \textbf{Causal inference}: Use propensity score matching or synthetic controls to estimate causal effects of injuries, rest, weather beyond correlation.
  \item \textbf{Probabilistic programming}: Implement full Bayesian models (Stan, Pyro) for better uncertainty quantification and prior elicitation.
\end{itemize}

\subsection{Operational Improvements}
\begin{itemize}
  \item \textbf{Private information}: Integrate injury monitoring (social media scraping, team beat reporters) to capture non-public intel before line moves.
  \item \textbf{Market microstructure arbitrage}: Exploit cross-book discrepancies in derivative markets (teaser pricing inconsistencies, correlated SGP legs).
  \item \textbf{Responsible gambling integration}: Build bankroll caps, session limits, and addiction detection into the system architecture.
  \item \textbf{Feature pipeline resilience}: The 2025 postseason re-evaluation (\S\ref{subsec:postseason-validation}) revealed that model prediction quality degrades severely when ancillary data sources (NextGen tracking, QBR, semantic features) become unavailable. Future production systems must implement fault-tolerant data pipelines with automated health monitoring and graceful degradation when individual data sources fail.
  \item \textbf{Postseason-specific modeling}: Playoff games ($n=13$ in 2025) showed distinct characteristics---tighter spreads, higher-quality opponents, motivational factors absent from regular-season training data. Dedicated playoff features (seed differential, bye-week advantage, playoff experience) merit investigation.
\end{itemize}

% Profit Optimization Future Work from Ch 12
% Profit Optimization Future Work
% Extracted from Chapter 12 for integration into Chapter 10

\subsection{Profit Optimization Extensions}\label{subsec:profit-optimization-extensions}

The profit optimization framework developed in this dissertation opens several research directions:

\begin{itemize}
  \item \textbf{Dynamic Edge Thresholds}: Current static thresholds could adapt based on recent performance and market conditions. Reinforcement learning could optimize threshold selection based on realized vs predicted edges, adjusting conservatism during hot/cold streaks.

  \item \textbf{Multi-Sport Extension}: Frameworks developed for NFL may transfer to other sports with adjustments. NBA and MLB offer higher game volumes for statistical learning, while international soccer markets may exhibit lower efficiency due to fragmented information flows.

  \item \textbf{Live Betting Integration}: In-game opportunities require real-time processing and different risk models. Bayesian updating of game state probabilities could exploit mispricing during momentum shifts, injury substitutions, and play-calling adjustments.

  \item \textbf{Automated Execution}: Current semi-automated approach could benefit from full automation with appropriate safeguards. API integration with sportsbooks, coupled with circuit breakers for anomalous line movements, would reduce latency and improve fill rates on profitable opportunities.

  \item \textbf{Market-Making Strategies}: Rather than purely price-taking, explore market-making on betting exchanges (Betfair, Pinnacle) to earn bid-ask spread while providing liquidity. This requires tighter spreads and faster computation but offers structural edge independent of directional predictions.

  \item \textbf{Correlation Arbitrage}: Exploit mispricing in correlated markets (team totals vs game total, first-half vs full-game spreads) using copula-based joint probability models. Same-game parlay pricing inconsistencies offer particularly attractive opportunities.
\end{itemize}

\begin{keyfinding}{Future Roadmap: Target Differential Efficiency}
The path to profitability runs through less-efficient market segments: live in-game betting (real-time Bayesian updates during momentum shifts), multi-sport transfer (NBA/MLB offer higher game volumes), and correlation arbitrage in same-game parlays where pricing inconsistencies persist.
\end{keyfinding}


\section{Broader Implications for Sports Analytics}

This dissertation clarifies the boundary between achievable and aspirational goals in sports betting---and demonstrates that \textbf{spread betting with public data falls in the ``not achievable'' category}:

\paragraph{Achieved (Methodology):}
\begin{itemize}
  \item Strong calibration (Brier $<$ 0.26 for NFL spreads)
  \item Positive CLV (beating closing lines on average, +14.9 basis points)
  \item \textbf{Rigorous leakage detection} (caught multiple target leakage issues across models)
  \item Conservative risk management (zero-bet weeks when OPE fails, 21\% abstention rate)
  \item Transparent methodology (reproducible pipelines, open-source code)
  \item \textbf{Honest negative result reporting} (51.0\% win rate, $-7.5$\% ROI after cleaning)
\end{itemize}

\paragraph{NOT Achieved (Economic):}
\begin{itemize}
  \item Systematic profitability (\textbf{NOT achieved}---cleaned baseline shows $-7.5$\% ROI)
  \item Win rate above breakeven (\textbf{NOT achieved}---51.0\% $<$ 52.4\% threshold)
  \item Positive Sharpe ratio (\textbf{NOT achieved}---Sharpe = $-1.22$ after cleaning)
\end{itemize}

\paragraph{Remaining Frontiers:}
\begin{itemize}
  \item \textbf{Less efficient markets}: Player props, live betting, alt lines may offer opportunities
  \item \textbf{Lower-vig exchanges}: Betfair, Pinnacle with 1--2\% vig instead of 4.5\%
  \item \textbf{Private information}: Insider injury intel, real-time tracking data
  \item Extension to live in-game markets (higher velocity, different information dynamics)
\end{itemize}

\textbf{The lesson for practitioners}: Prediction accuracy (74.9\%) does not imply betting profitability. The 4.5\% vigorish represents a substantial barrier that publicly available information cannot overcome. Future work should focus on less efficient markets or lower-friction execution venues.

\section{Final Reflection}

We began with optimism: perhaps rigorous methods could unlock NFL betting profits. We end with \textbf{clarity}: the comprehensive leakage investigation demonstrates that NFL spread betting markets are \textbf{efficient for publicly available information}. The cleaned baseline (51.0\% win rate, $-7.5$\% ROI) confirms this definitively.

The journey from v1.0 ($-83.5\%$ drawdown) through multiple iterations to the final cleaned baseline required:
\begin{itemize}
  \item Abandoning promising-looking approaches (BNN: 84 hours, zero production models)
  \item Discovering that simple principles matter more than complex algorithms (diversification $>$ sophistication)
  \item Catching catastrophic errors through multi-layer validation (Ridge Residuals target leakage: 88\% ``improvement'' $\to$ 47\% degradation)
  \item \textbf{Accepting negative results with the same rigor as positive ones}
\end{itemize}

The methods developed here---ensemble stacking with diversity optimization, copula-based dependence modeling, CVaR risk gates, OPE validation, and rigorous leakage detection---remain valuable for domains beyond NFL betting: portfolio optimization, resource allocation, and decision-making under uncertainty.

\vspace{1em}
\noindent \textbf{In summary}: This dissertation demonstrates that rigorous methods reveal truth, even when that truth is negative. Market efficiency is real. The contribution is methodological: we provide a template for how to rigorously validate (and invalidate) ML systems in adversarial environments, and how to honestly report negative results that advance the field.

\section{Closing Statement}

The methods developed here emphasize clarity and restraint over opacity and overfitting. We release code, evaluation scripts, and governance templates to lower barriers for future researchers studying market-facing systems under academic rigor.

The broader implication: AI systems deployed in efficient markets require explicit calibration, risk management, and transparent failure modes as first-class design goals. This work provides a template for such systems---and demonstrates that \textbf{honest negative results are as valuable as positive ones}. ``The market wins'' is not a failure of methodology; it is a confirmation of economic theory, rigorously validated.

\chaptersummary{
This dissertation investigated whether NFL betting markets could be beaten through disciplined methodology. The comprehensive data leakage investigation revealed that earlier positive ensemble results (pre-leakage +2.96\% ROI) were inflated by target leakage in the Ridge Residuals and CQL pipelines. The \textbf{cleaned baseline shows 51.0\% win rate and $-7.5$\% ROI}---confirming semi-strong market efficiency for spread betting with publicly available data. The simple power ratings model showed 57.7\% ATS in 2025, but multi-year data (51.8\% ATS, $-1.1$\% ROI) confirms it is not reliably profitable. The contribution is methodological: we demonstrate rigorous leakage detection across multi-season backtests, honest failure analysis, and transparent negative result reporting. The methods developed---ensemble diversification ($\bar{\rho} = 0.32$), copula-based dependence modeling, CVaR risk gates, OPE validation---remain valuable for domains beyond NFL betting.
}{
Future work should focus on extending the framework to \textbf{less efficient markets} (player props, live betting, lower-vig exchanges), implementing the five research vectors detailed in \Cref{sec:future-research} (causal features, transformers, LLMs, diversity optimization, MLOps), and validating any claimed edge on truly held-out data before deployment.
}
